\chapter{De bout en bout : la formulation quantique du hackathon}\label{ch:projet}
\iterbanner{V0}
\begin{evolution}{V1}{la boucle devient bi-objectif}Ce chapitre décrit la boucle de la V0, sur la seule traînée. En V1, elle a été étendue aux deux objectifs (traînée et SER), avec les contraintes codées dans le QUBO (13 qubits : réduction de Rosenberg et variables d'écart), et comparée à NSGA-II : voir le cours \emph{Optimisation bi-objectif traînée / SER}, chapitres 6 et 7.\end{evolution}

\section{La chaîne complète}
\begin{figure}[htbp]
\centering
\begin{tikzpicture}[node distance=5mm and 6mm]
\node[flow] (g) {Grille $8^3$\\ 9 bits};
\node[flow,right=of g] (s1) {Phase 1\\ aéro : $C_D$};
\node[flow,below=of s1] (s2) {Phase 2\\ SER};
\node[flow,right=12mm of s1,yshift=-5mm] (m) {Modèles\\ quadratiques\\ $\widehat{C_D}$, $\widehat{\rm SER}$};
\node[flow,right=of m,fill=corange!10] (qubo) {QUBO\\ $w\widehat{C_D}+(1-w)\widehat{\rm SER}$\\ (ch.~\ref{ch:qubo})};
\node[flow,right=of qubo,fill=cblue!10] (solve) {QAOA /\\ recuit\\ (ch.~\ref{ch:recuit}, \ref{ch:qaoa})};
\node[flow,below=12mm of qubo] (front) {front de Pareto\\ des profils \emph{simulés}};
\draw[arr] (g) -- (s1); \draw[arr] (g) |- (s2);
\draw[arr] (s1) -- (m); \draw[arr] (s2) -- (m);
\draw[arr] (m) -- (qubo); \draw[arr] (qubo) -- (solve);
\draw[arr] (solve) |- node[pos=0.25,right,font=\small] {candidats} (front);
\draw[arr] (front) -| node[pos=0.75,left,font=\small,align=right] {simuler les\\ nouveaux} (s2.south);
\end{tikzpicture}
\caption{La formulation quantique de la note d'intention. Le front final ne contient que des profils réellement simulés : le modèle quadratique ne sert qu'à \emph{proposer} des candidats.}\label{fig:chaine}
\end{figure}

\section{La boucle itérative, et pourquoi elle est indispensable}
\begin{tcolorbox}[coursbox,colback=white,colframe=cink,title={Algorithme (une seule fonction objectif, ici $C_D$).}]
\begin{enumerate}
\item Simuler $n_0$ profils (plan d'expériences initial).
\item Ajuster le modèle quadratique (10 coefficients) sur tous les profils simulés.
\item Construire le QUBO et le résoudre : QAOA, recuit, ou énumération quand $n$ est petit.
\item Simuler les $k$ meilleurs candidats \emph{non encore simulés}.
\item Recommencer jusqu'à stabilisation, ou épuisement du budget de simulations.
\end{enumerate}
\end{tcolorbox}
On a vu au chapitre~\ref{ch:qubo} qu'un modèle ajusté \emph{une fois} sur toute la base désigne un profil seulement 38\textsuperscript{e} sur 512. La boucle corrige ce défaut, car elle raffine le modèle là où sont les bons candidats.

\begin{figure}[htbp]
\centering
% figure générée par tools/make_figures.py : boucle_convergence
\begin{tikzpicture}
\begin{axis}[width=0.85\linewidth,height=5.6cm,axis lines=left,xmin=10,xmax=60,ymin=49,ymax=62,
  xlabel={nombre de profils simulés},ylabel={meilleur $C_D$ trouvé (counts), médiane},
  legend style={draw=none,font=\small,at={(0.98,0.98)},anchor=north east}]
\addplot[cblue,very thick] coordinates {(10,61) (11,60.92) (12,60.85) (13,60.79) (14,58.9) (15,58.51) (16,57) (17,54.78) (18,54.78) (19,53.41) (20,52.79) (21,52.62) (22,52.2) (23,51.16) (24,50.85) (25,50.53) (26,50.53) (27,50.53) (28,50.53) (29,50.53) (30,50.53) (31,50.53) (32,50.53) (33,50.53) (34,50.53) (35,50.53) (36,50.53) (37,50.53) (38,50.53) (39,50.53) (40,50.53) (41,50.53) (42,50.53) (43,50.53) (44,50.53) (45,50.53) (46,50.53) (47,50.53) (48,50.53) (49,50.53) (50,50.53) (51,50.53) (52,50.53) (53,50.53) (54,50.53) (55,50.53)};\addlegendentry{boucle : modèle quadratique, puis QUBO}
\addplot[corange,thick] coordinates {(10,61.1) (11,60) (12,58.9) (13,58.8) (14,58.65) (15,58.65) (16,58.51) (17,57.76) (18,57) (19,56.99) (20,56.99) (21,56.72) (22,55.95) (23,55.95) (24,55.16) (25,55.1) (26,55.1) (27,53.81) (28,53.81) (29,53.81) (30,53.81) (31,53.7) (32,53.7) (33,53.7) (34,53.06) (35,53.06) (36,53.03) (37,53.03) (38,53.03) (39,53.03) (40,53.03) (41,53.03) (42,53.03) (43,52.98) (44,52.79) (45,52.79) (46,52.79) (47,52.79) (48,52.79) (49,52.62) (50,52.62) (51,52.62) (52,52.62) (53,52.62) (54,52.62) (55,52.62) (56,52.62) (57,52.41) (58,52.41) (59,52.2) (60,52.2)};\addlegendentry{tirage aléatoire}
\addplot[cmuted,dashed] coordinates {(10,50.533) (60,50.533)};\addlegendentry{vrai minimum (50.5)}
\end{axis}
\end{tikzpicture}

\caption{Meilleur $C_D$ trouvé en fonction du nombre de profils simulés (médiane sur 50 tirages du plan initial : $n_0=10$ profils aléatoires, puis $k=3$ candidats par itération). Le QUBO est résolu ici par énumération, qui donne la même réponse qu'un QAOA ou un recuit parfaitement réussi. La boucle trouve le vrai minimum en \textbf{25 simulations} (médiane ; quartiles 22 et 34), contre 256 en moyenne pour un tirage au hasard sur 512 profils.}\label{fig:boucle}
\end{figure}

\begin{projet}[ : ce qu'il reste à construire]
Ce résultat ne porte que sur $C_D$. Il faut encore :
\begin{itemize}
\item ajouter la SER (phase~2) ;
\item traiter le bi-objectif (section suivante) ;
\item remplacer l'énumération par QAOA et le recuit ;
\item compter les simulations avec le même compteur que NSGA-II (\texttt{AeroSimulator}).
\end{itemize}
Les briques sont dans \texttt{Code/quantique/} : \texttt{qubo.py}, \texttt{qaoa.py}, \texttt{recuit.py}, \texttt{experiences.py}.
\end{projet}

\section{Deux objectifs, un seul QUBO}
Un QUBO n'a qu'une énergie. Pour obtenir un front de Pareto $(C_D,\ {\rm SER})$ :
\begin{itemize}
\item \textbf{Somme pondérée} : minimiser $w\,\tilde C_D+(1-w)\,\widetilde{\rm SER}$ (objectifs normalisés) pour plusieurs $w\in[0,1]$. Elle reste quadratique, donc c'est un QUBO, mais elle ne trouve que les points de la \emph{partie convexe} du front.
\item \textbf{$\varepsilon$-contrainte} : minimiser $\tilde C_D$ sous $\widetilde{\rm SER}\le\varepsilon$, avec des bits d'écart et une pénalité (ch.~\ref{ch:qubo}). Elle trouve tout le front, au prix de qubits supplémentaires, et la pénalité doit rester de degré 2.
\end{itemize}
\begin{figure}[htbp]
\centering
\begin{tikzpicture}
\begin{axis}[width=0.7\linewidth,height=5.6cm,axis lines=left,xmin=0,xmax=5,ymin=0,ymax=5,xtick=\empty,ytick=\empty,
  xlabel={$\tilde C_D$},ylabel={$\widetilde{\rm SER}$}]
\addplot[only marks,mark=*,cmuted,mark size=1.5pt] coordinates {(1.6,4.6)(2.4,4.1)(3.2,3.5)(2.7,2.9)(3.8,2.6)(4.3,1.9)(3.4,4.4)(4.6,3.6)(2.2,3.6)(4.0,4.0)};
\addplot[cblue,thick,mark=*,mark size=2.2pt] coordinates {(0.6,4.4)(1.2,2.4)(2.9,1.1)(4.4,0.5)};
\addplot[corange,only marks,mark=square*,mark size=2.6pt] coordinates {(2.3,2.0)};
\draw[cblue,dashed] (axis cs:1.2,2.4) -- (axis cs:2.9,1.1);
\node[corange,font=\small,anchor=west] at (axis cs:2.4,2.1) {non supporté};
\node[cblue,font=\small,anchor=west] at (axis cs:0.65,4.6) {front de Pareto};
\draw[-{Stealth},cmuted] (axis cs:4.4,4.6) -- (axis cs:3.9,4.1) node[pos=0,anchor=south,font=\small] {$w$ fixé : lignes de niveau};
\end{axis}
\end{tikzpicture}
\caption{Le point orange est Pareto-optimal (aucun autre profil n'est meilleur sur les deux objectifs), mais il est situé \emph{au-dessus} du segment qui relie ses voisins : aucune pondération $w$ ne le rend minimal. La somme pondérée ne le trouvera jamais.}\label{fig:pareto}
\end{figure}

\section{La comparaison avec NSGA-II}
Pour que la comparaison soit honnête :
\begin{itemize}
\item même grille de profils et même point de vol ;
\item même \textbf{compteur de simulations uniques} ;
\item même critère d'arrêt, fixé à l'avance : le front exact retrouvé, ou un pourcentage de son hypervolume ;
\item plusieurs graines aléatoires, avec médiane et dispersion.
\end{itemize}
Le front exact est connu par énumération des 512 profils. C'est ce qui rend la démonstration rigoureuse.

\section{Ressources et passage à l'échelle}
Avec $d$ variables de forme codées sur $k$ bits chacune, en encodage binaire :
\begin{center}
\begin{tabular}{lccc}
\toprule
& notre cas ($d=3$, $k=3$) & $d=10$, $k=4$ & $d=50$, $k=4$\\
\midrule
qubits $n=dk$ & 9 & 40 & 200\\
couplages (graphe complet) $\binom n2$ & 36 & 780 & 19\,900\\
CNOT par couche QAOA ($2\binom n2$) & 72 & 1\,560 & 39\,800\\
coefficients du modèle ($\binom{d+2}2$) & 10 & 66 & 1\,326\\
simulation par vecteur d'état & triviale & $\approx16$~To & impossible\\
\bottomrule
\end{tabular}
\end{center}
Deux enseignements à présenter :
\begin{itemize}
\item le modèle quadratique en variables physiques demande peu de coefficients, donc peu de simulations, alors qu'un QUBO libre sur les bits en demanderait $\binom n2+n+1$ ;
\item le graphe complet est le vrai obstacle matériel : profondeur QAOA en $O(n^2)$ par couche, et plongement coûteux sur un recuit quantique.
\end{itemize}

\section{Messages pour le jury}
\begin{enumerate}
\item La chaîne complète, de la physique au QUBO puis à QAOA et au recuit, est construite et \emph{vérifiée} sur un cas résolu exactement.
\item La qualité se mesure en \textbf{nombre de simulations coûteuses} pour retrouver le front exact, et non en temps de calcul quantique.
\item Le goulot d'étranglement est la \emph{fidélité du modèle quadratique}, pas la résolution du QUBO : la boucle itérative est essentielle.
\item Les ressources quantiques croissent en $O(n^2)$ portes par couche pour un QUBO dense. C'est l'argument à confronter au matériel actuel (NISQ).
\end{enumerate}

\exercices
\begin{exercice}
Trois profils ont pour objectifs normalisés $A=(0,\,1)$, $B=(0{,}6,\,0{,}6)$, $C=(1,\,0)$. (a) Sont-ils tous Pareto-optimaux ? (b) Montrer qu'aucun $w\in[0,1]$ ne fait de $B$ le minimum de $w\,x+(1-w)\,y$. (c) Quelle méthode permet de trouver $B$ ?
\end{exercice}
\begin{exercice}
Un cas industriel a 20 variables de forme à 16 niveaux chacune. Nombre de qubits en binaire ? Nombre de CNOT par couche de QAOA si le QUBO est dense ? Combien de simulations au minimum pour ajuster le modèle quadratique en variables physiques ? Et un QUBO libre sur les bits ?
\end{exercice}
